\clearpage
\section{CONSIDERAÇÃO FINAL}

\subsection{Análise Final}

O resultado central obtido até esta versão é a passagem de uma percepção genérica de que ``louças quebram'' para um problema de cadeia com fronteira, atores e consequências identificáveis. A formulação distingue o provável momento de geração do dano, o momento de detecção e o ator que absorve o custo. Essa separação reduz o risco de selecionar prematuramente uma solução que apenas facilite a conferência ou transfira a perda para outro participante da cadeia.

O recorte ainda contém incerteza material. A convergência de relatos orienta a investigação, mas não mede prevalência nem causalidade. O projeto informacional deverá reduzir essas incertezas antes da geração de concepções, distinguindo necessidades compartilhadas de necessidades específicas e verificando conflitos como maior proteção versus menor custo, menor volume e menor tempo operacional.

\subsection{Conclusão}

A etapa de pré-desenvolvimento delimitou um problema concreto: avarias de louças sanitárias associadas ao transporte entre o distribuidor e a obra, com danos que podem permanecer ocultos até etapas posteriores. No projeto informacional, a equipe adotou produção própria segundo a estratégia ATO e estruturou o ciclo operacional da solução desde o fornecedor até o descarte. O mapeamento cobre a entrega direta ao distribuidor, o uso em sua produção, o envio à obra, o recebimento e a movimentação pelo almoxarifado, a coleta de retorno e a destinação final. Os atores impactados foram classificados como clientes e receberam um conjunto preliminar de necessidades rastreado por fase. A próxima rodada deverá validar essas necessidades na linguagem dos próprios atores antes de convertê-las em requisitos, priorizá-las e estabelecer características mensuráveis para o produto.

\subsection{Lições Aprendidas}

\etapafutura{As lições aprendidas serão registradas ao longo das validações e consolidadas ao término do desenvolvimento.}
